
\documentclass{article}
\usepackage{colortbl}
\usepackage{makecell}
\usepackage{multirow}
\usepackage{supertabular}

\begin{document}

\newcounter{utterance}

\centering \large Interaction Transcript for game `function\_detective', experiment `logic\_passive\_examples', episode 193 with Qwen3.6{-}35B{-}A3B{-}FP8{-}without{-}reasoning.
\vspace{24pt}

{ \footnotesize  \setcounter{utterance}{1}
\setlength{\tabcolsep}{0pt}
\begin{supertabular}{c@{$\;$}|p{.15\linewidth}@{}p{.15\linewidth}p{.15\linewidth}p{.15\linewidth}p{.15\linewidth}p{.15\linewidth}}
    \# & \multicolumn{2}{c}{Player} && \multicolumn{2}{c}{Game Master} \\
    \hline

    \theutterance \stepcounter{utterance}  
    & & & \multicolumn{4}{p{0.6\linewidth}}{
        \cellcolor[rgb]{0.9,0.9,0.9}{
            \makecell[{{p{\linewidth}}}]{
                \texttt{\tiny{[P1$\langle$GM]}}
                \texttt{I have a hidden Python function for you to figure out.} \\
\\ 
\texttt{Your job is to infer what it does and then submit a correct Python implementation.} \\
\\ 
\texttt{Function signature:} \\
\texttt{(a: bool, b: bool) {-}> bool} \\
\\ 
\texttt{You have 15 turns total.} \\
\\ 
\texttt{I control which valid input{-}output examples you see. You can request them one at a time.} \\
\\ 
\texttt{You must make your best SOLVE attempt no later than the second{-}to{-}last turn.} \\
\texttt{Do not wait until the very end.} \\
\\ 
\texttt{You can use these actions:} \\
\texttt{{-} JSON format described at the end of the prompt} \\
\texttt{{-} SOLVE: ```python ... ```} \\
\\ 
\texttt{Do not add import statements in your guess.} \\
\texttt{Anything you might need is already available, so just write the function body in hypothesis and return the answer.} \\
\\ 
\texttt{Preloaded examples:} \\
\texttt{1. ((True, True), True)} \\
\\ 
\texttt{Scoring:} \\
\texttt{{-} correctness is required for any non{-}zero score} \\
\texttt{{-} in interactive modes, solving with fewer information requests is better} \\
\\ 
\texttt{Input formatting reminders:} \\
\texttt{{-} strings must be quoted} \\
\texttt{{-} booleans must be True/False} \\
\texttt{{-} use commas between multiple values} \\
\texttt{{-} do not write any extra action tags in the same message} \\
\\ 
\texttt{Format example:} \\
\texttt{You write:} \\
\texttt{\{'current\_mode': <probing/confirming/don't know; choose whichever best describes your current state in the task>, 'input\_rationale': <brief description of why you chose this hypothesis>, 'hypothesis': <either the exact string "unknown" if you have no guess yet OR your current best guess as runnable Python matching the given signature, with a one{-}line docstring stating what it does. Never write placeholder code, if you cannot commit to a guess, use "unknown". Match the given signature.} \\
\texttt{e.g.} \\
\texttt{def f(x):} \\
\texttt{	"""Returns the absolute value of x."""} \\
\texttt{	return x if x >= 0 else {-}x} \\
\\ 
\texttt{Use "unknown" if you have no guess.>, 'confidence\_score': <score between 1{-}100. 100 means highest confidence in your hypothesis>\}} \\
\\ 
\texttt{Output only one valid json per turn.} \\
            }
        }
    }
     \\ \\

    \theutterance \stepcounter{utterance}  
    & \multicolumn{4}{p{0.6\linewidth}}{
        \cellcolor[rgb]{0.561,0.737,0.561}{
            \makecell[{{p{\linewidth}}}]{
                \texttt{\tiny{[P1$\rangle$GM]}}
                \texttt{\{'current\_mode': 'probing', 'input\_rationale': 'I have one example: f(True, True) = True. This is consistent with several boolean functions (AND, OR, NAND, etc.). I need more data points to distinguish between them. I will request the output for (True, False) next.', 'hypothesis': 'unknown', 'confidence\_score': 10\}} \\
            }
        }
    }
    & & \\ \\

    \theutterance \stepcounter{utterance}  
    & & & \multicolumn{4}{p{0.6\linewidth}}{
        \cellcolor[rgb]{0.9,0.9,0.9}{
            \makecell[{{p{\linewidth}}}]{
                \texttt{\tiny{[P1$\langle$GM]}}
                \texttt{Game Violation: ParseError: Response is not valid JSON.} \\
\texttt{Output ONLY:} \\
\texttt{One JSON with confidence\_score, current\_mode, input\_rationale, and hypothesis keys ...} \\
\texttt{OR} \\
\texttt{SOLVE: ```python ... ```} \\
            }
        }
    }
     \\ \\

    \theutterance \stepcounter{utterance}  
    & \multicolumn{4}{p{0.6\linewidth}}{
        \cellcolor[rgb]{0.561,0.737,0.561}{
            \makecell[{{p{\linewidth}}}]{
                \texttt{\tiny{[P1$\rangle$GM]}}
                \texttt{\{} \\
\texttt{  "current\_mode": "probing",} \\
\texttt{  "input\_rationale": "I have one example: f(True, True) = True. This is consistent with several boolean functions (AND, OR, NAND, etc.). I need more data points to distinguish between them. I will request the output for (True, False) next.",} \\
\texttt{  "hypothesis": "unknown",} \\
\texttt{  "confidence\_score": 10} \\
\texttt{\}} \\
            }
        }
    }
    & & \\ \\

    \theutterance \stepcounter{utterance}  
    & & & \multicolumn{4}{p{0.6\linewidth}}{
        \cellcolor[rgb]{0.9,0.9,0.9}{
            \makecell[{{p{\linewidth}}}]{
                \texttt{\tiny{[P1$\langle$GM]}}
                \texttt{OUTPUT: ((True, False), False)} \\
            }
        }
    }
     \\ \\

    \theutterance \stepcounter{utterance}  
    & \multicolumn{4}{p{0.6\linewidth}}{
        \cellcolor[rgb]{0.561,0.737,0.561}{
            \makecell[{{p{\linewidth}}}]{
                \texttt{\tiny{[P1$\rangle$GM]}}
                \texttt{\{} \\
\texttt{  "current\_mode": "probing",} \\
\texttt{  "input\_rationale": "I now have two examples: f(True, True) = True and f(True, False) = False. This suggests the function might be equivalent to 'a' (identity on the first argument) or 'a AND b'. I need to check the case where a is False to distinguish between these and other possibilities like OR or NAND.",} \\
\texttt{  "hypothesis": "unknown",} \\
\texttt{  "confidence\_score": 50} \\
\texttt{\}} \\
            }
        }
    }
    & & \\ \\

    \theutterance \stepcounter{utterance}  
    & & & \multicolumn{4}{p{0.6\linewidth}}{
        \cellcolor[rgb]{0.9,0.9,0.9}{
            \makecell[{{p{\linewidth}}}]{
                \texttt{\tiny{[P1$\langle$GM]}}
                \texttt{OUTPUT: ((False, True), True)} \\
            }
        }
    }
     \\ \\

    \theutterance \stepcounter{utterance}  
    & \multicolumn{4}{p{0.6\linewidth}}{
        \cellcolor[rgb]{0.561,0.737,0.561}{
            \makecell[{{p{\linewidth}}}]{
                \texttt{\tiny{[P1$\rangle$GM]}}
                \texttt{\{} \\
\texttt{  "current\_mode": "probing",} \\
\texttt{  "input\_rationale": "I have three examples: f(True, True) = True, f(True, False) = False, f(False, True) = True. This eliminates AND (since f(False, True) is True, not False) and OR (since f(True, False) is False, not True). The pattern suggests the function returns True if b is True, regardless of a. This is consistent with the OR operation if we consider the inputs might be swapped or if it's simply 'b'. Let's check the last case f(False, False) to confirm if it's 'b' or something else like XOR (which would be False for (False, False) but True for (False, True) and (True, False) {-} wait, XOR is True for (True, False) which contradicts our second example). So XOR is out. The remaining strong candidate is 'b' (identity on second argument) or 'a OR b' if I misremembered OR. Let's re{-}evaluate: OR(1,0)=1, but we got 0. So not OR. AND(0,1)=0, but we got 1. So not AND. XOR(1,0)=1, but we got 0. So not XOR. XNOR(1,0)=0, XNOR(0,1)=0, but we} \\
            }
        }
    }
    & & \\ \\

    \theutterance \stepcounter{utterance}  
    & & & \multicolumn{4}{p{0.6\linewidth}}{
        \cellcolor[rgb]{0.9,0.9,0.9}{
            \makecell[{{p{\linewidth}}}]{
                \texttt{\tiny{[P1$\langle$GM]}}
                \texttt{Game Violation: ParseError: Response is not valid JSON.} \\
\texttt{Output ONLY:} \\
\texttt{One JSON with confidence\_score, current\_mode, input\_rationale, and hypothesis keys ...} \\
\texttt{OR} \\
\texttt{SOLVE: ```python ... ```} \\
            }
        }
    }
     \\ \\

    \theutterance \stepcounter{utterance}  
    & \multicolumn{4}{p{0.6\linewidth}}{
        \cellcolor[rgb]{0.561,0.737,0.561}{
            \makecell[{{p{\linewidth}}}]{
                \texttt{\tiny{[P1$\rangle$GM]}}
                \texttt{\{} \\
\texttt{  "current\_mode": "probing",} \\
\texttt{  "input\_rationale": "I have three examples: f(True, True) = True, f(True, False) = False, f(False, True) = True. This eliminates AND, OR, XOR, and XNOR. The pattern strongly suggests the function returns the value of the second argument 'b'. To confirm, I need to check the case where both are False.",} \\
\texttt{  "hypothesis": "unknown",} \\
\texttt{  "confidence\_score": 75} \\
\texttt{\}} \\
            }
        }
    }
    & & \\ \\

    \theutterance \stepcounter{utterance}  
    & & & \multicolumn{4}{p{0.6\linewidth}}{
        \cellcolor[rgb]{0.9,0.9,0.9}{
            \makecell[{{p{\linewidth}}}]{
                \texttt{\tiny{[P1$\langle$GM]}}
                \texttt{OUTPUT: ((False, False), True)} \\
            }
        }
    }
     \\ \\

    \theutterance \stepcounter{utterance}  
    & \multicolumn{4}{p{0.6\linewidth}}{
        \cellcolor[rgb]{0.561,0.737,0.561}{
            \makecell[{{p{\linewidth}}}]{
                \texttt{\tiny{[P1$\rangle$GM]}}
                \texttt{\{} \\
\texttt{  "current\_mode": "confirming",} \\
\texttt{  "input\_rationale": "I have all four input{-}output pairs: f(True, True)=True, f(True, False)=False, f(False, True)=True, f(False, False)=True. Let's analyze the truth table for 'b': b=True gives True, b=False gives False. This matches f(True, True)=True and f(True, False)=False. However, f(False, True)=True matches b=True, but f(False, False)=True does NOT match b=False (which would be False). So it's not simply 'b'. Let's look at 'a OR b': T|T=T, T|F=T (but we got F). No. 'a AND b': T\&T=T, T\&F=F, F\&T=F (but we got T). No. 'a XNOR b' (equivalence): T==T=T, T==F=F, F==T=F (but we got T). No. 'a XOR b': T\^{}T=F (but we got T). No. Let's look at the outputs: T, F, T, T. Inputs: TT, TF, FT, FF. Output is False only when input is TF. This is the NAND of (a AND NOT b)? No. It is equivalent to NOT (a AND NOT b). Which is a {-}> b (implication). Let's check} \\
            }
        }
    }
    & & \\ \\

    \theutterance \stepcounter{utterance}  
    & & & \multicolumn{4}{p{0.6\linewidth}}{
        \cellcolor[rgb]{0.9,0.9,0.9}{
            \makecell[{{p{\linewidth}}}]{
                \texttt{\tiny{[P1$\langle$GM]}}
                \texttt{Game Violation: ParseError: Response is not valid JSON.} \\
\texttt{Output ONLY:} \\
\texttt{One JSON with confidence\_score, current\_mode, input\_rationale, and hypothesis keys ...} \\
\texttt{OR} \\
\texttt{SOLVE: ```python ... ```} \\
            }
        }
    }
     \\ \\

    \theutterance \stepcounter{utterance}  
    & \multicolumn{4}{p{0.6\linewidth}}{
        \cellcolor[rgb]{0.561,0.737,0.561}{
            \makecell[{{p{\linewidth}}}]{
                \texttt{\tiny{[P1$\rangle$GM]}}
                \texttt{\{} \\
\texttt{  "current\_mode": "confirming",} \\
\texttt{  "input\_rationale": "I have all four input{-}output pairs: f(True, True)=True, f(True, False)=False, f(False, True)=True, f(False, False)=True. The only case where the output is False is when a=True and b=False. This corresponds to the logical implication a {-}> b, which is equivalent to (not a) or b.",} \\
\texttt{  "hypothesis": "def f(a: bool, b: bool) {-}> bool:\textbackslash{}n    \textbackslash{}"\textbackslash{}"\textbackslash{}"Returns the logical implication a {-}> b, which is equivalent to (not a) or b.\textbackslash{}"\textbackslash{}"\textbackslash{}"\textbackslash{}n    return (not a) or b",} \\
\texttt{  "confidence\_score": 100} \\
\texttt{\}} \\
            }
        }
    }
    & & \\ \\

    \theutterance \stepcounter{utterance}  
    & & & \multicolumn{4}{p{0.6\linewidth}}{
        \cellcolor[rgb]{0.9,0.9,0.9}{
            \makecell[{{p{\linewidth}}}]{
                \texttt{\tiny{[P1$\langle$GM]}}
                \texttt{OUTPUT: ((True, True), True)} \\
            }
        }
    }
     \\ \\

    \theutterance \stepcounter{utterance}  
    & \multicolumn{4}{p{0.6\linewidth}}{
        \cellcolor[rgb]{0.561,0.737,0.561}{
            \makecell[{{p{\linewidth}}}]{
                \texttt{\tiny{[P1$\rangle$GM]}}
                \texttt{\{} \\
\texttt{  "current\_mode": "confirming",} \\
\texttt{  "input\_rationale": "I have all four input{-}output pairs: f(True, True)=True, f(True, False)=False, f(False, True)=True, f(False, False)=True. The only case where the output is False is when a=True and b=False. This corresponds to the logical implication a {-}> b, which is equivalent to (not a) or b. The new example ((True, True), True) is consistent with this hypothesis.",} \\
\texttt{  "hypothesis": "def f(a: bool, b: bool) {-}> bool:\textbackslash{}n    \textbackslash{}"\textbackslash{}"\textbackslash{}"Returns the logical implication a {-}> b, which is equivalent to (not a) or b.\textbackslash{}"\textbackslash{}"\textbackslash{}"\textbackslash{}n    return (not a) or b",} \\
\texttt{  "confidence\_score": 100} \\
\texttt{\}} \\
            }
        }
    }
    & & \\ \\

    \theutterance \stepcounter{utterance}  
    & & & \multicolumn{4}{p{0.6\linewidth}}{
        \cellcolor[rgb]{0.9,0.9,0.9}{
            \makecell[{{p{\linewidth}}}]{
                \texttt{\tiny{[P1$\langle$GM]}}
                \texttt{OUTPUT: ((True, False), False)} \\
            }
        }
    }
     \\ \\

    \theutterance \stepcounter{utterance}  
    & \multicolumn{4}{p{0.6\linewidth}}{
        \cellcolor[rgb]{0.561,0.737,0.561}{
            \makecell[{{p{\linewidth}}}]{
                \texttt{\tiny{[P1$\rangle$GM]}}
                \texttt{\{} \\
\texttt{  "current\_mode": "confirming",} \\
\texttt{  "input\_rationale": "I have all four input{-}output pairs: f(True, True)=True, f(True, False)=False, f(False, True)=True, f(False, False)=True. The only case where the output is False is when a=True and b=False. This corresponds to the logical implication a {-}> b, which is equivalent to (not a) or b. The new example ((True, False), False) is consistent with this hypothesis.",} \\
\texttt{  "hypothesis": "def f(a: bool, b: bool) {-}> bool:\textbackslash{}n    \textbackslash{}"\textbackslash{}"\textbackslash{}"Returns the logical implication a {-}> b, which is equivalent to (not a) or b.\textbackslash{}"\textbackslash{}"\textbackslash{}"\textbackslash{}n    return (not a) or b",} \\
\texttt{  "confidence\_score": 100} \\
\texttt{\}} \\
            }
        }
    }
    & & \\ \\

    \theutterance \stepcounter{utterance}  
    & & & \multicolumn{4}{p{0.6\linewidth}}{
        \cellcolor[rgb]{0.9,0.9,0.9}{
            \makecell[{{p{\linewidth}}}]{
                \texttt{\tiny{[P1$\langle$GM]}}
                \texttt{OUTPUT: ((False, True), True)} \\
            }
        }
    }
     \\ \\

    \theutterance \stepcounter{utterance}  
    & \multicolumn{4}{p{0.6\linewidth}}{
        \cellcolor[rgb]{0.561,0.737,0.561}{
            \makecell[{{p{\linewidth}}}]{
                \texttt{\tiny{[P1$\rangle$GM]}}
                \texttt{\{} \\
\texttt{  "current\_mode": "confirming",} \\
\texttt{  "input\_rationale": "I have all four input{-}output pairs: f(True, True)=True, f(True, False)=False, f(False, True)=True, f(False, False)=True. The only case where the output is False is when a=True and b=False. This corresponds to the logical implication a {-}> b, which is equivalent to (not a) or b. The new example ((False, True), True) is consistent with this hypothesis.",} \\
\texttt{  "hypothesis": "def f(a: bool, b: bool) {-}> bool:\textbackslash{}n    \textbackslash{}"\textbackslash{}"\textbackslash{}"Returns the logical implication a {-}> b, which is equivalent to (not a) or b.\textbackslash{}"\textbackslash{}"\textbackslash{}"\textbackslash{}n    return (not a) or b",} \\
\texttt{  "confidence\_score": 100} \\
\texttt{\}} \\
            }
        }
    }
    & & \\ \\

    \theutterance \stepcounter{utterance}  
    & & & \multicolumn{4}{p{0.6\linewidth}}{
        \cellcolor[rgb]{0.9,0.9,0.9}{
            \makecell[{{p{\linewidth}}}]{
                \texttt{\tiny{[P1$\langle$GM]}}
                \texttt{OUTPUT: ((False, False), True)} \\
            }
        }
    }
     \\ \\

    \theutterance \stepcounter{utterance}  
    & \multicolumn{4}{p{0.6\linewidth}}{
        \cellcolor[rgb]{0.561,0.737,0.561}{
            \makecell[{{p{\linewidth}}}]{
                \texttt{\tiny{[P1$\rangle$GM]}}
                \texttt{\{} \\
\texttt{  "current\_mode": "confirming",} \\
\texttt{  "input\_rationale": "I have all four input{-}output pairs: f(True, True)=True, f(True, False)=False, f(False, True)=True, f(False, False)=True. The only case where the output is False is when a=True and b=False. This corresponds to the logical implication a {-}> b, which is equivalent to (not a) or b. The new example ((False, False), True) is consistent with this hypothesis. I am confident in this solution.",} \\
\texttt{  "hypothesis": "def f(a: bool, b: bool) {-}> bool:\textbackslash{}n    \textbackslash{}"\textbackslash{}"\textbackslash{}"Returns the logical implication a {-}> b, which is equivalent to (not a) or b.\textbackslash{}"\textbackslash{}"\textbackslash{}"\textbackslash{}n    return (not a) or b",} \\
\texttt{  "confidence\_score": 100} \\
\texttt{\}} \\
            }
        }
    }
    & & \\ \\

    \theutterance \stepcounter{utterance}  
    & & & \multicolumn{4}{p{0.6\linewidth}}{
        \cellcolor[rgb]{0.9,0.9,0.9}{
            \makecell[{{p{\linewidth}}}]{
                \texttt{\tiny{[P1$\langle$GM]}}
                \texttt{OUTPUT: ((True, True), True)} \\
            }
        }
    }
     \\ \\

    \theutterance \stepcounter{utterance}  
    & \multicolumn{4}{p{0.6\linewidth}}{
        \cellcolor[rgb]{0.561,0.737,0.561}{
            \makecell[{{p{\linewidth}}}]{
                \texttt{\tiny{[P1$\rangle$GM]}}
                \texttt{\{} \\
\texttt{  "current\_mode": "confirming",} \\
\texttt{  "input\_rationale": "I have verified all four possible input combinations for the function f(a, b). The outputs are: f(True, True)=True, f(True, False)=False, f(False, True)=True, f(False, False)=True. This truth table uniquely identifies the logical implication a {-}> b, which is logically equivalent to (not a) or b. I am confident in this solution.",} \\
\texttt{  "hypothesis": "def f(a: bool, b: bool) {-}> bool:\textbackslash{}n    \textbackslash{}"\textbackslash{}"\textbackslash{}"Returns the logical implication a {-}> b, which is equivalent to (not a) or b.\textbackslash{}"\textbackslash{}"\textbackslash{}"\textbackslash{}n    return (not a) or b",} \\
\texttt{  "confidence\_score": 100} \\
\texttt{\}} \\
            }
        }
    }
    & & \\ \\

    \theutterance \stepcounter{utterance}  
    & & & \multicolumn{4}{p{0.6\linewidth}}{
        \cellcolor[rgb]{0.9,0.9,0.9}{
            \makecell[{{p{\linewidth}}}]{
                \texttt{\tiny{[P1$\langle$GM]}}
                \texttt{OUTPUT: ((True, False), False)} \\
            }
        }
    }
     \\ \\

    \theutterance \stepcounter{utterance}  
    & \multicolumn{4}{p{0.6\linewidth}}{
        \cellcolor[rgb]{0.561,0.737,0.561}{
            \makecell[{{p{\linewidth}}}]{
                \texttt{\tiny{[P1$\rangle$GM]}}
                \texttt{\{} \\
\texttt{  "current\_mode": "confirming",} \\
\texttt{  "input\_rationale": "I have verified all four possible input combinations for the function f(a, b). The outputs are: f(True, True)=True, f(True, False)=False, f(False, True)=True, f(False, False)=True. This truth table uniquely identifies the logical implication a {-}> b, which is logically equivalent to (not a) or b. I am confident in this solution.",} \\
\texttt{  "hypothesis": "def f(a: bool, b: bool) {-}> bool:\textbackslash{}n    \textbackslash{}"\textbackslash{}"\textbackslash{}"Returns the logical implication a {-}> b, which is equivalent to (not a) or b.\textbackslash{}"\textbackslash{}"\textbackslash{}"\textbackslash{}n    return (not a) or b",} \\
\texttt{  "confidence\_score": 100} \\
\texttt{\}} \\
            }
        }
    }
    & & \\ \\

    \theutterance \stepcounter{utterance}  
    & & & \multicolumn{4}{p{0.6\linewidth}}{
        \cellcolor[rgb]{0.9,0.9,0.9}{
            \makecell[{{p{\linewidth}}}]{
                \texttt{\tiny{[P1$\langle$GM]}}
                \texttt{OUTPUT: ((False, True), True)} \\
            }
        }
    }
     \\ \\

    \theutterance \stepcounter{utterance}  
    & \multicolumn{4}{p{0.6\linewidth}}{
        \cellcolor[rgb]{0.561,0.737,0.561}{
            \makecell[{{p{\linewidth}}}]{
                \texttt{\tiny{[P1$\rangle$GM]}}
                \texttt{\{} \\
\texttt{  "current\_mode": "confirming",} \\
\texttt{  "input\_rationale": "I have verified all four possible input combinations for the function f(a, b). The outputs are: f(True, True)=True, f(True, False)=False, f(False, True)=True, f(False, False)=True. This truth table uniquely identifies the logical implication a {-}> b, which is logically equivalent to (not a) or b. I am confident in this solution.",} \\
\texttt{  "hypothesis": "def f(a: bool, b: bool) {-}> bool:\textbackslash{}n    \textbackslash{}"\textbackslash{}"\textbackslash{}"Returns the logical implication a {-}> b, which is equivalent to (not a) or b.\textbackslash{}"\textbackslash{}"\textbackslash{}"\textbackslash{}n    return (not a) or b",} \\
\texttt{  "confidence\_score": 100} \\
\texttt{\}} \\
            }
        }
    }
    & & \\ \\

    \theutterance \stepcounter{utterance}  
    & & & \multicolumn{4}{p{0.6\linewidth}}{
        \cellcolor[rgb]{0.9,0.9,0.9}{
            \makecell[{{p{\linewidth}}}]{
                \texttt{\tiny{[P1$\langle$GM]}}
                \texttt{OUTPUT: ((False, False), True)} \\
            }
        }
    }
     \\ \\

    \theutterance \stepcounter{utterance}  
    & \multicolumn{4}{p{0.6\linewidth}}{
        \cellcolor[rgb]{0.561,0.737,0.561}{
            \makecell[{{p{\linewidth}}}]{
                \texttt{\tiny{[P1$\rangle$GM]}}
                \texttt{SOLVE: ```python} \\
\texttt{def f(a: bool, b: bool) {-}> bool:} \\
\texttt{    """Returns the logical implication a {-}> b, which is equivalent to (not a) or b."""} \\
\texttt{    return (not a) or b} \\
\texttt{```} \\
            }
        }
    }
    & & \\ \\

    \theutterance \stepcounter{utterance}  
    & & & \multicolumn{2}{p{0.3\linewidth}}{
        \cellcolor[rgb]{0.9,0.9,0.9}{
            \makecell[{{p{\linewidth}}}]{
                \texttt{\tiny{[GM$|$GM]}}
                \texttt{Game Verdict: WIN} \\
            }
        }
    }
    & & \\ \\

    \theutterance \stepcounter{utterance}  
    & & & \multicolumn{2}{p{0.3\linewidth}}{
        \cellcolor[rgb]{0.9,0.9,0.9}{
            \makecell[{{p{\linewidth}}}]{
                \texttt{\tiny{[GM$|$GM]}}
                \texttt{That is correct!} \\
            }
        }
    }
    & & \\ \\

    \theutterance \stepcounter{utterance}  
    & & & \multicolumn{2}{p{0.3\linewidth}}{
        \cellcolor[rgb]{0.9,0.9,0.9}{
            \makecell[{{p{\linewidth}}}]{
                \texttt{\tiny{[GM$|$GM]}}
                \texttt{def a\_requires\_b(a: bool, b: bool) {-}> bool:} \\
\texttt{    if a and not b:} \\
\texttt{        return False} \\
\texttt{    return True} \\
            }
        }
    }
    & & \\ \\

    \theutterance \stepcounter{utterance}  
    & & & \multicolumn{2}{p{0.3\linewidth}}{
        \cellcolor[rgb]{0.9,0.9,0.9}{
            \makecell[{{p{\linewidth}}}]{
                \texttt{\tiny{[GM$|$GM]}}
                \texttt{def f(a: bool, b: bool) {-}> bool:} \\
\texttt{    """Returns the logical implication a {-}> b, which is equivalent to (not a) or b."""} \\
\texttt{    return (not a) or b} \\
            }
        }
    }
    & & \\ \\

    \theutterance \stepcounter{utterance}  
    & & & \multicolumn{2}{p{0.3\linewidth}}{
        \cellcolor[rgb]{0.9,0.9,0.9}{
            \makecell[{{p{\linewidth}}}]{
                \texttt{\tiny{[GM$|$GM]}}
                \texttt{Mode: passive\_examples} \\
\texttt{Accuracy: 1.000} \\
\texttt{Information requests: 11/15} \\
\texttt{Revealed passive examples: 12} \\
\texttt{Efficiency: 0.333} \\
\texttt{Slack: 4} \\
\texttt{Observed{-}pair consistency: 1.000} \\
\texttt{Parse errors: 3} \\
\texttt{Runtime errors: 0} \\
            }
        }
    }
    & & \\ \\

\end{supertabular}
}

\end{document}
